\chapter{Implementing a Simple Interpreter  \label{chapter:interpreter}}
Figure \ref{fig:Pure.g} on page \pageref{fig:Pure.g} shows the grammar of a simple programming language.
In this chapter, we will develop an \blue{interpreter}\index{interpreter} for this language.  This is our most ambitious
example:  In contrast to the calculator of Section \ref{section:calculator} and the differentiator of Section
\ref{section:differentiator}, the language now has \blue{statements} and hence a \blue{control flow}.

Figure \ref{fig:sum.sl} shows
an example program that conforms to this grammar. This program first reads a number
which is stored in the variable \texttt{n}. Subsequently, the sum
\\[0.2cm]
\hspace*{1.3cm}
$\ds s := \sum\limits_{\mbox{$\normalsize i=1$}}^{\mbox{$\normalsize n^2$}} \mbox{\Large $i$}$
\\[0.2cm]
is computed using a \texttt{while} loop.  Finally, this sum is printed.

\begin{figure}[!ht]
\centering
\begin{Verbatim}[ frame         = lines, 
                  framesep      = 0.3cm, 
                  labelposition = bottomline,
                  numbers       = left,
                  numbersep     = -0.2cm,
                  xleftmargin   = 0.8cm,
                  xrightmargin  = 0.8cm,
                ]
    n := read();   // read a number
    s := 0;
    i := 0;
    while (i < n * n) {
        i := i + 1;
        s := s + i;
    }
    /* print the result */
    print(s);
\end{Verbatim}
\vspace*{-0.3cm}
\caption{A program to compute the sum $\sum\limits_{i=1}^{n^2} i$.}
\label{fig:sum.sl}
\end{figure}

\begin{figure}[!ht]

  \begin{center}
  \framebox{
  \framebox{
  \begin{minipage}[t]{12.4cm}
  \begin{eqnarray*}
  \textsl{program}      & \rightarrow & \;\varepsilon                                                              \\
                        & \mid        & \;\textsl{stmnt}\; \textsl{program}                                        \\[0.2cm]
  \textsl{stmnt}        & \rightarrow &   \quoted{if} \quoted{(} \textsl{bool\_expr} \quoted{)}
                                          \quoted{\{} \textsl{program} \quoted{\}}                                 \\
                        & \mid        &   \quoted{while} \quoted{(} \textsl{bool\_expr} \quoted{)}
                                          \quoted{\{} \textsl{program} \quoted{\}}                                 \\
                        & \mid        & \;\simtextsc{Id} \quoted{:=} \textsl{expr} \quoted{;}                      \\
                        & \mid        & \;\textsl{expr} \quoted{;}                                                 \\[0.2cm]
  \textsl{bool\_expr}   & \rightarrow & \;\textsl{expr} \quoted{==} \textsl{expr}                                  \\
                        & \mid        & \;\textsl{expr} \quoted{!=} \textsl{expr}                                  \\
                        & \mid        & \;\textsl{expr} \quoted{<=} \textsl{expr}                                  \\
                        & \mid        & \;\textsl{expr} \quoted{>=} \textsl{expr}                                  \\
                        & \mid        & \;\textsl{expr} \quoted{<}  \textsl{expr}                                  \\
                        & \mid        & \;\textsl{expr} \quoted{>}  \textsl{expr}                                  \\[0.2cm]
  \textsl{expr}         & \rightarrow & \;\textsl{expr} \quoted{+} \textsl{product}                                \\
                        & \mid        & \;\textsl{expr} \quoted{-} \textsl{product}                                \\
                        & \mid        & \;\textsl{product}                                                         \\[0.2cm]
  \textsl{product}      & \rightarrow & \;\textsl{product} \quoted{*} \textsl{factor}                              \\
                        & \mid        & \;\textsl{product} \quoted{/} \textsl{factor}                              \\
                        & \mid        & \;\textsl{product} \quoted{\%} \textsl{factor}                             \\
                        & \mid        & \;\textsl{factor}                                                          \\[0.2cm]
  \textsl{factor}       & \rightarrow &   \quoted{(} \textsl{expr} \quoted{)}                                      \\
                        & \mid        & \;\simtextsc{Int}                                                          \\
                        & \mid        & \;\simtextsc{Id}                                                           \\
                        & \mid        & \;\simtextsc{Id} \quoted{(} \textsl{expr\_list} \quoted{)}                 \\[0.2cm]
  \textsl{expr\_list}   & \rightarrow & \;\varepsilon                                                              \\
                        & \mid        & \;\textsl{ne\_expr\_list}                                                  \\[0.2cm]
  \textsl{ne\_expr\_list} & \rightarrow & \;\textsl{expr}                                                          \\
                        & \mid        & \;\textsl{expr} \quoted{,} \textsl{ne\_expr\_list}
  \end{eqnarray*}
  \vspace*{-0.5cm}
  \end{minipage}}}
  \end{center}
  \caption{A grammar for a simple programming language.}
  \label{fig:Pure.g}
\end{figure}

The statements of the language that we want to implement include \textsl{assignments}, \textsl{function calls},
\textsl{\texttt{if} statements}, and \textsl{\texttt{while} loops}.  The language supports boolean expressions
that can be built with the help of the comparison operators
\\[0.2cm]
\hspace*{1.3cm}
``\texttt{==}'', ``\texttt{!=}'', ``\texttt{<=}'', ``\texttt{>=}'', ``\texttt{<}'', and ``\texttt{>}''.
\\[0.2cm]
Arithmetic expressions can use the arithmetic operators
\\[0.2cm]
\hspace*{1.3cm}
``\texttt{+}'', ``\texttt{-}'', ``\texttt{*}'', ``\texttt{/}'', and ``\texttt{\%}''.
\\[0.2cm]
In the base case, arithmetic expressions are built from numbers, variables, and function calls.
Although the language supports the calling of predefined functions like \texttt{read} and \texttt{print},
it does not support user defined functions.

In contrast to the language \texttt{C}, the body of an \texttt{if} statement or a \texttt{while} loop always
has to be enclosed in braces, even if it consists of a single statement.  Hence, the body is a
\textsl{program}, i.e.~a possibly empty sequence of statements.  A pair of braces can only occur as part of an
\texttt{if} statement or a \texttt{while} loop.  If the language is extended by \texttt{else} branches, this rule
ensures that the dangling else problem, which is discussed in Chapter \ref{chapter:compiler}, cannot
arise.


Similar to our program for symbolic differentiation, we represent the individual statements as
nested tuples.  The program shown in Figure \ref{fig:sum.sl} is represented by the nested tuple displayed in
Figure \ref{fig:sum.ast}.  This nested tuple is nothing other than the abstract syntax tree
corresponding to the original program.  Note that lists of statements are represented as nested tuples
that start with the string \squoted{program}.

\begin{figure}[!ht]
\centering
\begin{Verbatim}[ frame         = lines,
                  framesep      = 0.3cm,
                  labelposition = bottomline,
                  numbers       = left,
                  numbersep     = 0.3cm,
                  xleftmargin   = 0.3cm,
                  xrightmargin  = 0.0cm,
                  gobble        = 4,
                ]
    ('program',
        (':=', ('ID', 'n'), ('call', ('ID', 'read'))),
        (':=', ('ID', 's'), ('INT', '0')),
        (':=', ('ID', 'i'), ('INT', '0')),
        ('while',
            ('<', ('ID', 'i'), ('*', ('ID', 'n'), ('ID', 'n'))),
            ('program',
                (':=', ('ID', 'i'), ('+', ('ID', 'i'), ('INT', '1'))),
                (':=', ('ID', 's'), ('+', ('ID', 's'), ('ID', 'i')))
            )
        ),
        ('expr', ('call', ('ID', 'print'), ('ID', 's')))
    )
\end{Verbatim}
\vspace*{-0.3cm}
\caption{A nested tuple that represents the program from Figure \ref{fig:sum.sl}.}
\label{fig:sum.ast}
\end{figure}
\FloatBarrier

\noindent
The figures \ref{fig:01-Interpreter.ipynb:grammar-1} and \ref{fig:01-Interpreter.ipynb:grammar-2} on page
\pageref{fig:01-Interpreter.ipynb:grammar-1} show the transcription of the grammar of Figure
\ref{fig:Pure.g} into the \simtextsc{Ebnf} dialect of \simtextsc{Larkita}.  Recall that there is no separate
scanner specification:  Both the lexical structure and the syntactical structure of our language are
described by the single string \texttt{sl\_grammar}.  Figure \ref{fig:01-Interpreter.ipynb:grammar-1}
contains the \blue{syntactical} part of this string, while Figure \ref{fig:01-Interpreter.ipynb:grammar-2}
contains its \blue{lexical} part.

\begin{figure}[!ht]
\centering
\begin{minted}[ frame        = lines,
                framesep     = 0.3cm,
                bgcolor      = sepia,
                numbers      = left,
                numbersep    = 0.3cm,
                xleftmargin  = 0.3cm,
                xrightmargin = 0.0cm,
              ]{python3}
    from larkita import Larkita

    sl_grammar = r"""
    program: stmnt*

    stmnt: "if"    "(" bool_expr ")" "{" program "}"
         | "while" "(" bool_expr ")" "{" program "}"
         | ID ":=" expr ";"
         | expr ";"                                 -> "expr"

    bool_expr: expr ("==" | "!=" | "<=" | ">=" | "<" | ">") expr

    expr: expr ("+" | "-") product
        | product

    product: product ("*" | "/" | "%") factor
           | factor

    factor: "(" expr ")"
          | INT
          | ID
          | ID "(" [expr ("," expr)*] ")" -> "call"

    // These literals are left out of the heads of all tuples.
    %ignore: "(", ")", "{", "}", ";"
\end{minted}
\vspace*{-0.3cm}
\caption{The grammar of our programming language, part \texttt{I}: syntax.}
\label{fig:01-Interpreter.ipynb:grammar-1}
\end{figure}

\begin{figure}[!ht]
\centering
\begin{minted}[ frame        = lines,
                framesep     = 0.3cm,
                bgcolor      = sepia,
                firstnumber  = 27,
                numbers      = left,
                numbersep    = 0.3cm,
                xleftmargin  = 0.3cm,
                xrightmargin = 0.0cm,
              ]{text}
    // Lexical definitions
    %import common.CNAME -> ID
    %import common.INT
    %import common.WS

    // In order to show how it's done, I have defined the tokens
    // CPP_COMMENT and C_COMMENT describing comments myself.
    CPP_COMMENT: /\/\/[^\n]*/
    C_COMMENT:   "/*" /(.|\n)*?/ "*/"

    // Ignore whitespace and comments
    %ignore WS, C_COMMENT, CPP_COMMENT
    """

    parser = Larkita(sl_grammar, out_file='sl.out')
\end{minted}
\vspace*{-0.3cm}
\caption{The grammar of our programming language, part \texttt{II}: lexical definitions.}
\label{fig:01-Interpreter.ipynb:grammar-2}
\end{figure}

\noindent
As the general structure of this grammar specification is already familiar from the previous two sections,
we only discuss those parts that are new.
\begin{enumerate}
\item Line 3 opens the string \texttt{sl\_grammar}.  Note the prefix ``\texttt{r}'':  This string is a
      \blue{raw string}.  In a raw string, the backslash loses its special meaning, i.e.~it does not start
      an \blue{escape sequence} but rather denotes itself.  This is necessary because the regular
      expression in line 34 contains the character sequence ``\verb$\/$'', %$
      which is \emph{not} a valid
      escape sequence of \textsl{Python}.  Inside an ordinary string, this sequence would provoke a
      warning.
\item Line 4 defines the variable \texttt{program}, which is the start symbol, since it is the first variable
      of the grammar.  Instead of transcribing the two grammar rules
      \\[0.2cm]
      \hspace*{1.3cm}
      $\textsl{program} \rightarrow \varepsilon$
      \quad and \quad
      $\textsl{program} \rightarrow \textsl{stmnt}\;\textsl{program}$
      \\[0.2cm]
      literally, we use the \blue{Kleene star} and write \texttt{stmnt*}.  Here, the
      \simtextsc{Ebnf} notation of \simtextsc{Larkita} really pays off:  The abstract syntax trees of
      \emph{all} the statements that have been found become the children of a single tuple.  Had we used the
      recursive formulation instead, the abstract syntax tree would have been a deeply nested chain of tuples.

      As the alternative \texttt{stmnt*} contains no literal, the name of the rule is used as the head of
      the tuple.  Hence, a program consisting of the statements $s_1$, $\cdots$, $s_n$ is represented as the tuple
      \texttt{('program', $s_1$, $\cdots$, $s_n$)}.  Note that the head of a tuple is determined by the grammar
      and not by the input:  Even a program that consists of a single statement is represented by a tuple
      with the head \squoted{program}.
\item Line 6 -- 9 define the variable \texttt{stmnt}.  The four alternatives correspond one to one to the
      four alternatives of the grammar shown in Figure \ref{fig:Pure.g}.  The keywords \texttt{"if"} and
      \texttt{"while"} are written as literals, i.e.~in exactly the same way as the operator symbols.  In
      particular, we do \emph{not} have to declare tokens for them.
\item Line 25 is a variant of the directive \texttt{\%ignore} that is written with a colon:  The directive
      \\[0.2cm]
      \hspace*{1.3cm}
      \verb|%ignore: "(", ")", "{", "}", ";"|
      \\[0.2cm]
      removes the given literals from the heads of \emph{all} alternatives of the grammar, since the
      parentheses, the braces, and the semicolon carry no information.  Therefore, the alternative in line 6
      yields a tuple of the form \texttt{('if', $\textsl{test}$, $\textsl{body}$)}, where $\textsl{body}$ is
      the abstract syntax tree of the \texttt{program} in braces, while the assignment in line 8 yields a tuple
      of the form
      \\[0.2cm]
      \hspace*{1.3cm}
      \texttt{(':=', ('ID', $\textsl{name}$), $\textsl{expr}$)}.
      \\[0.2cm]
      As the body of an \texttt{if} statement or a \texttt{while} loop is a \texttt{program}, it is
      represented by a tuple that starts with the string \squoted{program}, even if it consists of a single
      statement.  This is the reason why the body of the \texttt{while} loop in Figure \ref{fig:sum.ast} is a
      tuple that starts with the string \squoted{program}.  The alternative \texttt{expr ";"} in line 9
      consists of a single symbol, since the semicolon is ignored.  Without the alias \texttt{"expr"}, it would return the abstract syntax tree
      of the expression and we could not distinguish an expression statement from an expression.  The alias
      replaces the head, so an expression statement is represented as \texttt{('expr', $\textsl{expr}$)}.
      Note that an alias may be written with or without double quotes.
\item Line 11 defines the variable \texttt{bool\_expr}.  The part
      \\[0.2cm]
      \hspace*{1.3cm}
      \verb+("==" | "!=" | "<=" | ">=" | "<" | ">")+
      \\[0.2cm]
      is an inner alternative.  As it does not occur inside a repetition, it is an abbreviation (see
      Section \ref{section:calculator-ebnf}):  Line 11 stands for
      the six alternatives \texttt{expr "==" expr}, $\cdots$, \texttt{expr ">" expr}, one for every
      comparison operator.  Therefore, the comparison operator that has been read is the head of the
      resulting tuple, e.g.~a test of the form $\textsl{lhs}$ \texttt{==} $\textsl{rhs}$ yields the tuple
      \texttt{('==', $\textsl{lhs}$, $\textsl{rhs}$)}.
\item Line 13 -- 17 define the variables \texttt{expr} and \texttt{product}.  These rules encode the
      \blue{precedence} of the arithmetic operators in the usual way.  Again, inner alternatives are used:
      Line 13 stands for the two alternatives \texttt{expr "+" product} and \texttt{expr "-" product}.  In
      addition to the four basic arithmetic operators, our language supports the
      \blue{modulo operator} \squoted{\%}.
\item Line 19 -- 22 define the variable \texttt{factor}.  As the parentheses are ignored, the alternative in
      line 19 returns the abstract syntax tree of the expression in parentheses.  Line 22 describes a
      \blue{function call}.  The list of arguments is written as
      \\[0.2cm]
      \hspace*{1.3cm}
      \texttt{[expr ("," expr)*]}.
      \\[0.2cm]
      The square brackets make the part enclosed in them \blue{optional}, while the Kleene star describes an
      arbitrary number of repetitions.  Therefore, this single line replaces the two grammar rules for
      \textsl{expr\_list} and the two grammar rules for \textsl{ne\_expr\_list} that are shown in Figure
      \ref{fig:Pure.g}.  The children of an optional part and of a repetition are spliced into the tuple
      of the alternative and the alias \texttt{"call"} replaces the head.  Therefore, a function call of the
      form $f(a_1, \cdots\!, a_n)$ is represented as the tuple
      \\[0.2cm]
      \hspace*{1.3cm}
      \texttt{('call', ('ID', $f$), $a_1$, $\cdots$, $a_n$)}.
\item Line 28 -- 30 import three token definitions from the library \texttt{common}.  In contrast to the
      calculator of Section \ref{section:calculator}, we rename the token \texttt{CNAME} to the shorter name
      \texttt{ID}, and we import the token \texttt{INT} instead of the token \texttt{NUMBER}, because the
      language that we implement only supports integers.
\item Line 34 -- 35 define the tokens \texttt{CPP\_COMMENT} and \texttt{C\_COMMENT}.  In \simtextsc{Larkita},
      a \blue{regular expression} is written between two slashes, exactly as in the language
      \textsl{JavaScript}.  Our language supports two kinds of comments.
      \begin{enumerate}
      \item Line 34 describes \blue{one line comments} that start with the character sequence
            ``\texttt{//}'' and extend to the end of the line.  Note that the two slashes that start the
            comment have to be escaped inside the regular expression, since an unescaped slash would
            terminate it.  This is the reason why \texttt{sl\_grammar} has to be a raw string.
      \item Line 35 describes \blue{multiline comments}, i.e.~text that is surrounded by ``\texttt{/*}''
            and ``\texttt{*/}''.

            Here, we have to use the so-called \emph{non-greedy} version of the Kleene operator
            ``\texttt{*}''.  The non-greedy version of the operator ``\texttt{*}'' is written as
            ``\texttt{*?}'' and matches as little as possible.  Therefore, the regular expression
            \\[0.2cm]
            \hspace*{1.3cm}
            \verb$"/*" /(.|\n)*?/ "*/"$ %$
            \\[0.2cm]
            represents a string that starts with the character sequence ``\texttt{/*}'', ends with the
            character sequence ``\texttt{*/}'', and is as short as possible.  This ensures that in a line
            like
            \\[0.2cm]
            \hspace*{1.3cm}
            \texttt{/* Hugo */ i := i + 1; /* Anton */}
            \\[0.2cm]
            two separate comments are recognized.
      \end{enumerate}
\item Line 38 instructs the scanner to discard white space and both kinds of comments.  The terminals that
      are to be ignored are listed after \texttt{\%ignore}, separated by commas.
\item Line 41 creates the LR(1) parser.  A description of the states of this parser is written into the file
      \texttt{sl.out}.
\end{enumerate}
\FloatBarrier

\noindent
The following table summarizes the shapes of the nested tuples that are returned by the parser.
\begin{center}
\begin{tabular}{|l|l|}
\hline
\textbf{construct}     & \textbf{nested tuple} \\
\hline
program, body          & \texttt{('program', $\textsl{stmnt}_1$, $\cdots$, $\textsl{stmnt}_n$)} \\
assignment             & \texttt{(':=', ('ID', $\textsl{name}$), $\textsl{expr}$)} \\
expression statement   & \texttt{('expr', $\textsl{expr}$)} \\
conditional, loop      & \texttt{('if', $\textsl{test}$, $\textsl{body}$)}, \texttt{('while', $\textsl{test}$, $\textsl{body}$)} \\
binary operator        & \texttt{($\textsl{op}$, $\textsl{lhs}$, $\textsl{rhs}$)}, for example \texttt{('+', $\textsl{lhs}$, $\textsl{rhs}$)} \\
number, variable       & \texttt{('INT', $\textsl{digits}$)}, \texttt{('ID', $\textsl{name}$)} \\
function call          & \texttt{('call', ('ID', $\textsl{name}$), $\textsl{arg}_1$, $\cdots$, $\textsl{arg}_n$)} \\
\hline
\end{tabular}
\end{center}

\exampleEng
The string
\\[0.2cm]
\hspace*{1.3cm}
\texttt{"x := f(1, 2+3, y);"}
\\[0.2cm]
is parsed into the nested tuple
\\[0.2cm]
\hspace*{0.3cm}
\texttt{('program', (':=', ('ID', 'x'),}
\\
\hspace*{0.8cm}
\texttt{('call', ('ID', 'f'), ('INT', '1'),}
\\
\hspace*{1.3cm}
\texttt{('+', ('INT', '2'), ('INT', '3')), ('ID', 'y'))))}.
\hspace*{\fill} $\diamond$

\remarkEng
The treatment of the keywords ``\texttt{if}'' and ``\texttt{while}'' deserves a closer look.  Both of these
strings also match the token \texttt{ID}, which has been imported as \blue{\texttt{common.CNAME}}.  A
scanner that is generated with the module \texttt{re} would therefore have to recognize these keywords as
identifiers first and then look them up in a dictionary of keywords in order to change the type of the
token.  With \simtextsc{Larkita}, none of this is necessary:  Tokens that are given as \blue{literal strings}
automatically have a higher priority than tokens that are given as regular expressions.  Consequently, the
string ``\texttt{while}'' is recognized as the keyword, while the string ``\texttt{whilepig}'' is
recognized as an identifier.  \eox

\remarkEng
The library \texttt{common} already provides the definitions \texttt{C\_COMMENT} and \texttt{CPP\_COMMENT}.
Hence, line 34 -- 35 of Figure \ref{fig:01-Interpreter.ipynb:grammar-2} could be replaced with two
\texttt{\%import} directives.  We have written these tokens ourselves in order to show how a
regular expression is specified in \simtextsc{Larkita}.  \eox
\vspace*{0.2cm}

\noindent
Figure \ref{fig:Interpreter.ipynb:main} shows the function \texttt{main}, which receives the name of a
file containing a program in our simple programming language as its argument.

\begin{figure}[!ht]
\centering
\begin{minted}[ frame        = lines,
                framesep     = 0.3cm,
                bgcolor      = sepia,
                numbers      = left,
                numbersep    = 0.3cm,
                xleftmargin  = 0.8cm,
                xrightmargin = 0.8cm,
              ]{python3}
    def main(file):
        with open(file, 'r') as handle:
            program = handle.read()
        ast = parser.parse(program)
        display(tuple2dot(ast))
        Values = {}
        execute(ast, Values)
\end{minted}
\vspace*{-0.3cm}
\caption{Specification of the function \texttt{main}.}
\label{fig:Interpreter.ipynb:main}
\end{figure}

\noindent
The function \texttt{main} performs the steps that are by now familiar.
\begin{enumerate}
\item Line 2 -- 3 read the program into the string \texttt{program}.
\item Line 4 parses this string.  The result is the abstract syntax tree of the program, which is represented
      as a nested tuple.
\item Line 5 displays this abstract syntax tree graphically.  The function \texttt{tuple2dot} is defined in
      the auxiliary notebook
      \href{https://github.com/karlstroetmann/Formal-Languages/blob/master/Python/Chapter-08/AST2Dot.ipynb}{\texttt{AST2Dot.ipynb}},
      which renders a nested tuple with the help of the
      \href{https://graphviz.org}{\textsl{Graphviz}} library.  This function has nothing to do with
      parsing, but it is very useful when debugging a grammar:  It makes the structure of the abstract
      syntax tree immediately visible.

      The notebook also provides the function \texttt{parse}, which reads and parses a program and displays
      both the program text and the resulting nested tuple.  This function is handy if we want to inspect an
      abstract syntax tree without executing the corresponding program.
\item Line 6 -- 7 execute the abstract syntax tree.  In order to do so, the function \texttt{execute} needs
      to know the values of all variables.  These are stored in the dictionary \texttt{Values}, which
      initially is empty.  Every time a variable is assigned, the variable and the corresponding value are
      stored in this dictionary.
\end{enumerate}

\begin{figure}[!ht]
\centering
\begin{minted}[ frame        = lines,
                framesep     = 0.3cm,
                bgcolor      = sepia,
                numbers      = left,
                numbersep    = 0.3cm,
                xleftmargin  = 0.3cm,
                xrightmargin = 0.0cm,
              ]{python3}
    def execute(stmnt, Values):
        match stmnt:
            case ('program', *SL):
                execute_tuple(tuple(SL), Values)
            case (':=', ('ID', var), value):
                Values[var] = evaluate(value, Values)
            case ('expr', expr):
                evaluate(expr, Values)
            case ('if', test, stmnt):
                if evaluate_bool(test, Values):
                    execute(stmnt, Values)
            case ('while', test, stmnt):
                while evaluate_bool(test, Values):
                    execute(stmnt, Values)
            case _:
                assert False, f'{stmnt} unexpected'

    def execute_tuple(StatementList, Values={}):
        for stmnt in StatementList:
            execute(stmnt, Values)
\end{minted}
\vspace*{-0.3cm}
\caption{The function \texttt{execute}.}
\label{fig:Interpreter.ipynb:execute}
\end{figure}

Figure \ref{fig:Interpreter.ipynb:execute} shows the implementation of the function \texttt{execute}.  This
implementation consists mainly of a large case distinction based on the type of the statement to be
executed.  Note that this function does not refer to \simtextsc{Larkita} in any way:  It operates on the pure
\textsl{Python} object that has been returned by the parser.
\begin{enumerate}
\item First, we check if \texttt{stmnt} is a list of statements.  A statement is a list of statements
      if the first element of the tuple representing this statement is the string \squoted{program}.

      In this case, the statements following the string \squoted{program} are executed via the auxiliary
      function \texttt{execute\_tuple}.  Observe that the pattern in line 3 uses a
      \blue{star pattern}\index{pattern matching!star pattern}:  The variable \texttt{SL} is bound to the list of
      all remaining components of the tuple.  This is necessary because a list of statements can have any
      length.
\item If the statement is an assignment of the form
      \\[0.2cm]
      \hspace*{1.3cm}
      \texttt{(':=', ('ID', var), value)}
      \\[0.2cm]
      then the value of the expression \texttt{value} is computed using the function \texttt{evaluate}.
      This value is then stored in the dictionary \texttt{Values} under the key \texttt{var}.  Note that the
      pattern in line 5 extracts the name of the variable from the token that represents it.
\item If the statement is an expression statement of the form
      \\[0.2cm]
      \hspace*{1.3cm}
      \texttt{('expr', expr)}
      \\[0.2cm]
      then the corresponding \texttt{expr} is evaluated using the function \texttt{evaluate}.
\item If the statement is of the form
      \\[0.2cm]
      \hspace*{1.3cm}
      \texttt{('if', test, stmnt)}
      \\[0.2cm]
      then \texttt{test} is a nested tuple representing a Boolean expression, while \texttt{stmnt} is the body,
      i.e.~a list of statements.
      In this case, the expression \texttt{test} is first evaluated using the function \texttt{evaluate\_bool}.
      If this evaluation yields the value \texttt{True}, then \texttt{stmnt} is executed.
\item If the statement is of the form
      \\[0.2cm]
      \hspace*{1.3cm}
      \texttt{('while', test, stmnt)}
      \\[0.2cm]
      then \texttt{test} is a Boolean expression and \texttt{stmnt} is the body of the loop.
      In this case, the expression \texttt{test} is first evaluated using the function \texttt{evaluate\_bool}.
      If this evaluation yields the value \texttt{True}, then \texttt{stmnt} is executed.
      Then the expression \texttt{test} is evaluated again.  If the result is
      \texttt{False}, then the execution of the while loop finishes.  Otherwise,
      \texttt{stmnt} is executed as long as the evaluation of \texttt{test}
      yields \texttt{True}.
\end{enumerate}

\begin{figure}[!ht]
\centering
\begin{minted}[ frame        = lines, 
                framesep     = 0.3cm, 
                bgcolor      = sepia,
                numbers      = left,
                numbersep    = 0.4cm,
                xleftmargin  = 0.0cm,
                xrightmargin = 0.0cm,
              ]{python3}
def evaluate_bool(expr: NestedTuple, Values: dict[str, Number]) -> bool:
    match expr:
        case ('==', lhs, rhs):
            return evaluate(lhs, Values) == evaluate(rhs, Values)
        case ('!=', lhs, rhs):
            return evaluate(lhs, Values) != evaluate(rhs, Values)
        case ('<=', lhs, rhs):
            return evaluate(lhs, Values) <= evaluate(rhs, Values)
        case ('>=', lhs, rhs):
            return evaluate(lhs, Values) >= evaluate(rhs, Values)
        case ('<', lhs, rhs):
            return evaluate(lhs, Values) <  evaluate(rhs, Values)
        case ('>', lhs, rhs):
            return evaluate(lhs, Values) >  evaluate(rhs, Values)
        case _:
            assert False, f'{expr} unexpected'
\end{minted}
\vspace*{-0.3cm}
\caption{The evaluation of boolean expressions.}
\label{fig:Interpreter.ipynb:evaluate_bool}
\end{figure}

Figure \ref{fig:Interpreter.ipynb:evaluate_bool} shows the implementation of the function \texttt{evaluate\_bool}.
This function receives a boolean expression and a dictionary storing the current values of the variables as input.
\begin{enumerate}[(a)]
\item If the expression to be evaluated is a Boolean expression of the form
      \\[0.2cm]
      \hspace*{1.3cm}
      \texttt{('==', lhs, rhs)}
      \\[0.2cm]
      we recursively evaluate the expressions \texttt{lhs} and \texttt{rhs} and return \texttt{True} if and only if
      both expressions yield the same value.
\item If the expression to be evaluated is a Boolean expression of the form
      \\[0.2cm]
      \hspace*{1.3cm}
      \texttt{('<', lhs, rhs)}
      \\[0.2cm]
      we recursively evaluate the expressions \texttt{lhs} and \texttt{rhs} and return \texttt{True} if and only if
      the value obtained from the evaluation of \texttt{lhs} is less than the
      value obtained from the evaluation of \texttt{rhs}.

      The remaining cases are similar.
\end{enumerate}

\begin{figure}[!ht]
\centering
\begin{minted}[ frame        = lines,
                framesep     = 0.3cm,
                bgcolor      = sepia,
                numbers      = left,
                numbersep    = 0.4cm,
                xleftmargin  = 0.0cm,
                xrightmargin = 0.0cm,
              ]{python3}
def evaluate(expr: NestedTuple, Values: dict[str, Number]) -> Number:
    match expr:
        case ('INT', digits):
            return int(digits)
        case ('ID', name):
            return Values[name]
        case ('call', ('ID', 'read')):
            return int(input('Please enter a natural number: '))
        case ('call', ('ID', 'print'), expr):
            print(evaluate(expr, Values))
            return 0
        case ('+', lhs, rhs):
            return evaluate(lhs, Values) + evaluate(rhs, Values)
        case ('-', lhs, rhs):
            return evaluate(lhs, Values) - evaluate(rhs, Values)
        case ('*', lhs, rhs):
            return evaluate(lhs, Values) * evaluate(rhs, Values)
        case ('/', lhs, rhs):
            return evaluate(lhs, Values) / evaluate(rhs, Values)
        case ('%', lhs, rhs):
            return evaluate(lhs, Values) % evaluate(rhs, Values)
        case _:
            assert False, f'{expr} unexpected'
\end{minted}
\vspace*{-0.3cm}
\caption{The evaluation of arithmetic expressions.}
\label{fig:Interpreter.ipynb:evaluate}
\end{figure}

Figure \ref{fig:Interpreter.ipynb:evaluate} shows the implementation of the function \texttt{evaluate}.
This function receives an arithmetic expression and a dictionary storing the
values of the variables as input.  The types \texttt{NestedTuple} and \texttt{Number} are defined in the
notebook:  \texttt{NestedTuple} is the type of the abstract syntax trees, while \texttt{Number} is the type of
the values of our variables.
\begin{enumerate}[(a)]
\item If the expression to be evaluated is a number, i.e.~a token of the form
      \\[0.2cm]
      \hspace*{1.3cm}
      \texttt{('INT', digits)},
      \\[0.2cm]
      we convert the string \texttt{digits} into an integer and return this integer as the result.
\item If the expression to be evaluated is a variable, i.e.~a token of the form
      \\[0.2cm]
      \hspace*{1.3cm}
      \texttt{('ID', name)},
      \\[0.2cm]
      we look up the value of this variable in the dictionary \texttt{Values} and return this value as the
      result.
\item If the expression to be evaluated is a call to the function \texttt{read},
      we prompt the user to enter a natural number.  We then convert the string entered by the user into an integer.

      Note that the pattern in line 7 is a tuple of length two:  An empty argument list produces no
      children at all and therefore the tuple representing the call \texttt{read()} is
      \\[0.2cm]
      \hspace*{1.3cm}
      \texttt{('call', ('ID', 'read'))}.
\item If the expression to be evaluated is a call to the function \texttt{print}, we evaluate the argument and
      print it.  Furthermore, we return 0 since the evaluation of an expression always has to return a value.
\item If the expression to be evaluated is a sum of the form
      \\[0.2cm]
      \hspace*{1.3cm}
      \texttt{('+', lhs, rhs)}
      \\[0.2cm]
      we recursively evaluate the expressions \texttt{lhs} and \texttt{rhs} and return the sum of these
      values. 
\item The evaluation of the arithmetic operators \texttt{'-'}, \texttt{'*'}, \texttt{'/'}, and \texttt{'\%'}
      is analogous to the evaluation of the operator \texttt{'+'}.
\end{enumerate}
\FloatBarrier

\noindent
The Jupyter notebook \texttt{01-Interpreter}, which is available at
\\[0.2cm]
\hspace*{1.3cm}
\href{https://github.com/karlstroetmann/Formal-Languages/blob/master/Python/Chapter-08/01-Interpreter.ipynb}{Python/Chapter-08/01-Interpreter.ipynb}
\\[0.2cm]
implements the interpreter that has been discussed in this chapter.

\exerciseEng
\begin{enumerate}[(a)]
\item Add \texttt{for} loops to the interpreter.
\item Expand the interpreter to include logical operators
      ``\texttt{\&\&}'' for logical \emph{and}, ``\texttt{||}'' for logical \emph{or},
      and ``\texttt{!}'' for \emph{negation}.  The operator ``\texttt{!}'' should bind 
      strongest and the operator ``\texttt{||}'' should bind weakest.
\item Enhance the interpreter to allow for user-defined functions.
      \begin{itemize}
      \item A function should only have access to its parameters and those variables that
            are defined locally inside the function.  
      \item A function should always return a value using a \texttt{return} statement.
            In order to facilitate that a function can contain multiple \texttt{return}
            statements and that a \texttt{return} statement can occur
            anywhere in the function, use an \blue{exception} to communicate the
            return value and to transfer the control flow out of the function.
      \end{itemize}
\end{enumerate}
Figure \ref{fig:factorial.sl} on page \pageref{fig:factorial.sl} shows a simple program
that implements a function to compute the factorial function.  Your interpreter should be able to execute this
program.
Use the Jupyter notebook \texttt{Interpreter-Frame}, which is available at
\\[0.2cm]
\hspace*{1.3cm}
\href{https://github.com/karlstroetmann/Formal-Languages/blob/master/Python/Chapter-08/Interpreter-Frame.ipynb}{Python/Chapter-08/Interpreter-Frame.ipynb}
\\[0.2cm]
as a starting point for your solution.
\\[0.2cm]
\hint
Both the \texttt{for} loop and the logical operators can be added by extending the grammar shown in Figure
\ref{fig:01-Interpreter.ipynb:grammar-1} and by adding the corresponding cases to the functions
\texttt{execute} and \texttt{evaluate\_bool}.  Make sure that the heads of the tuples that are built for the
new alternatives are the ones that your interpreter expects.  If necessary, give the new alternatives
aliases.  For the logical operators, introduce the syntactical variables \texttt{disjunction},
\texttt{conjunction}, and \texttt{negation} in order to encode the precedences.  As for \texttt{if} statements and
\texttt{while} loops, the bodies of \texttt{for} loops and of functions have to be enclosed in braces.
\eox

\begin{figure}[!ht]
\centering
\begin{minted}[ frame         = lines, 
                 framesep      = 0.3cm, 
                 firstnumber   = 1,
                 bgcolor       = sepia,
                 numbers       = left,
                 numbersep     = 0.3cm,
                 xleftmargin   = 0.8cm,
                 xrightmargin  = 0.8cm,
                 ]{javascript}
    function factorial(n) {
        if (n == 0) {
            return 1;
        }
        return n * factorial(n - 1);
    }
    
    for (i := 2; i <= 25; i := i + 1) {
        print(factorial(i));
    }                 
\end{minted}
\vspace*{-0.3cm}
\caption{A program to compute the factorial function.}
\label{fig:factorial.sl}
\end{figure}

\section{Check your Understanding}
\begin{enumerate}[(a)]
\item How can a grammar rule of the form
      \\[0.2cm]
      \hspace*{1.3cm}
      $\textsl{v} \rightarrow \varepsilon \;\mid\; \textsl{w}\; \textsl{v}$
      \\[0.2cm]
      be expressed more concisely in the \simtextsc{Ebnf} dialect of \simtextsc{Larkita}, and what does the
      resulting abstract syntax tree look like?
\item Why is the string that contains the grammar of our interpreter written as a
      \blue{raw string}?
\item What is the effect of the directive \verb|%ignore: "(", ")", "{", "}", ";"| in the grammar of our
      interpreter?
\item Why do the keywords of a programming language not have to be declared as tokens when the scanner is
      generated by \simtextsc{Larkita}?  How would this problem be solved if the scanner were implemented with
      the module \texttt{re}?
\item Why do we have to use the \emph{non-greedy} version of the Kleene operator when specifying multiline
      comments?
\item Are you able to extend the interpreter that has been discussed in this chapter?
\item Assume that we develop an interpreter that supports user defined functions.
      How can we accommodate the fact that a \texttt{return} statement can occur anywhere in a function?
\end{enumerate}


%%% Local Variables: 
%%% mode: latex
%%% TeX-master: "formal-languages.tex"
%%% End: 
